\chapter{ניהול סיכונים}\label{ch:risk}

ניהול הסיכונים נעשה לפי \en{ISO 31000}~\cite{iso31000}. לכל סיכון נקבעו הסתברות (\en{P}) והשפעה (\en{I}) בסולם \num{1} עד \num{5}, והחשיפה היא המכפלה של שתיהן. סיכון שחשיפתו \num{12} ומעלה מוגדר ״אדום״, ומנוהל בישיבת צוות שבועית. רשימת הסיכונים מתעדכנת בכל אבן דרך. לכל סיכון נקבע גם מועד סגירה: אבן הדרך שבה הוא צפוי לרדת מהרשימה, או ״שוטף״ אם הוא מלווה את המוצר לאורך כל חייו.

\small

ראו טבלה~\ref{tab:risks}.

\begin{longtable}{@{}l p{3.6cm} c c c p{4.4cm} c@{}}
\caption{רשימת סיכונים, חשיפה ותגובה}\label{tab:risks}\\
\toprule
\textbf{מזהה} & \textbf{סיכון} & \textbf{\en{P}} & \textbf{\en{I}} & \textbf{חשיפה} & \textbf{מניעה ותוכנית גיבוי} & \textbf{סגירה} \\
\midrule
\endfirsthead
\toprule
\textbf{מזהה} & \textbf{סיכון} & \textbf{\en{P}} & \textbf{\en{I}} & \textbf{חשיפה} & \textbf{מניעה ותוכנית גיבוי} & \textbf{סגירה} \\
\midrule
\endhead
\bottomrule
\endfoot
\req{R-01} & מחולל חדש עוקף את הגלאים & \num{5} & \num{4} & \textbf{20} & אימון חוזר רבעוני; הפרדה לפי מחולל; אזור ״לא ודאי״ מעביר לאדם. גיבוי: הרחבה זמנית של אזור הביניים & שוטף \\
\req{R-02} & לא מושג יעד שיעור התרעות שווא מעל \num{2\%} (פוגע באמון העורכים) & \num{3} & \num{5} & \textbf{15} & כיול על נתוני הלקוח; ״מסלול ירוק״. גיבוי: הזזת הסף העליון ותיעוד בבקשת שינוי & \en{M4} \\
\req{R-03} & עורכים עוקפים את התהליך בלחץ זמן & \num{4} & \num{4} & \textbf{16} & חסימה טכנית (\req{FR-14}); כלל \num{7} בתקן; מדד זמן החלטה. גיבוי: הדרכה חוזרת ודיווח להנהלה & \en{M6} \\
\req{R-04} & עיכוב ברישוי צילומים לאימון & \num{3} & \num{3} & \num{9} & חוזים עם \num{3} סוכנויות במקביל; התחלה מוקדמת בחודש \num{1} & \en{M1} \\
\req{R-05} & תעודות \en{C2PA} מנוצלות לרעה (מפתח שנגנב, פרצה במצלמה) & \num{2} & \num{5} & \num{10} & בדיקת רשימות ביטול בכל אימות; תעודה היא ראיה חזקה ולא הוכחה מוחלטת & שוטף \\
\req{R-06} & שינוי בהנחיות סעיף \num{50} של האיחוד האירופי & \num{3} & \num{3} & \num{9} & מעקב משפטי רבעוני; ארכיטקטורה מודולרית לתוויות & \en{M6} \\
\req{R-07} & דליפת תמונות לא מפורסמות (חומר עיתונאי רגיש) & \num{2} & \num{5} & \num{10} & הצפנה, מחיקה אחרי \num{30} יום, בדיקות חדירה, אחסון באזור הלקוח & שוטף \\
\req{R-08} & מערכות הפיילוט נוטשות באמצע & \num{2} & \num{4} & \num{8} & הסכם פיילוט חתום; \num{7} מועמדות ל-\num{5} מקומות & \en{M5} \\
\req{R-09} & חריגה בתקציב המחשוב & \num{3} & \num{2} & \num{6} & מכסות ענן; גלאי קל ראשון; רזרבה של \num{12\%} & \en{M4} \\
\req{R-10} & ״דיבידנד השקרן״: המערכת משמשת לפסילת תמונה אמיתית & \num{2} & \num{4} & \num{8} & תווית ״לא ודאי״ אינה ״מזויף״; כלל \num{6}; פרסום מתודולוגיה & \en{M6} \\
\end{longtable}
\normalsize

ראו איור~\ref{fig:heat}.


\begin{figure}[H]
\begin{ltrdiagram}
% heat: TikZ / pgfplots source, built inside ltrdiagram (LTR)
\begin{tikzpicture}[font=\sffamily\scriptsize]
\foreach \p in {1,...,5}{\foreach \i in {1,...,5}{
  \pgfmathsetmacro{\e}{\p*\i}
  \ifdim\e pt<6pt \def\cc{okc!25}\else\ifdim\e pt<12pt \def\cc{midc!35}\else \def\cc{badc!30}\fi\fi
  \fill[\cc,draw=white,line width=1.2pt] (\i*1.25-1.25,\p*0.95-0.95) rectangle ++(1.25,0.95);
}}
\foreach \i in {1,...,5}{\node at (\i*1.25-0.625,-0.3) {\i};}
\foreach \p in {1,...,5}{\node at (-0.3,\p*0.95-0.475) {\p};}
\node at (3.125,-0.75) {\he{השפעה \en{I}}};
\node[rotate=90] at (-0.85,2.375) {\he{הסתברות \en{P}}};
\foreach \r/\p/\i/\dx in {01/5/4/0, 02/3/5/0, 03/4/4/0, 04/3/3/-0.32, 06/3/3/0.32, 05/2/5/-0.32, 07/2/5/0.32, 08/2/4/-0.32, 10/2/4/0.32, 09/3/2/0}{
  \node[circle,fill=ink,text=white,inner sep=1.2pt,font=\sffamily\scriptsize\bfseries] at (\i*1.25-0.625+\dx,\p*0.95-0.475) {R\r};}
\end{tikzpicture}

\end{ltrdiagram}
\caption{מפת חום של הסיכונים (\num{5}$\times$\num{5}): שלושה סיכונים באזור האדום}\label{fig:heat}
\end{figure}

\textbf{תובנה מרכזית:} שני הסיכונים החמורים ביותר (\req{R-01} ו-\req{R-03}) אינם טכניים במובן הצר. האחד הוא מרוץ חימוש מול מחוללים, והשני הוא התנהגות אנושית בלחץ. לכן התגובה לשניהם משלבת טכנולוגיה עם נוהל: אזור ״לא ודאי״ וכלל \num{7} בתקן.
